% DOCX paragraph 40
\chapter*{El comienzo del viaje}\label{introduction}
\addcontentsline{toc}{chapter}{El comienzo del viaje}
\markboth{El comienzo del viaje}{El comienzo del viaje}

% DOCX paragraph 43
La gravedad atrae: es un fenómeno tan evidente que no hacía falta escribir este libro para subrayarlo. Menos evidente resulta que, incluso antes de manifestarse como interacción física, la gravedad atrae nuestra atención y nuestra imaginación. Desde que nacemos, antes de desarrollar una relación consciente con el universo físico, ya conocemos la gravedad de manera instintiva. Durante el resto de nuestra vida será la única de las cuatro interacciones fundamentales de la que tendremos conciencia. Y de la que a menudo intentaremos escapar.

% DOCX paragraph 44
He escrito este libro con el deseo de explicar qué es la gravedad y por qué nos sentimos irresistiblemente atraídos por ella, aunque sea de manera subconsciente. Para ello me propongo llevar a los lectores de viaje por sus efectos observables y, en particular, por esas regiones de la física que nos ha revelado la revolucionaria teoría de la gravitación de Einstein: la relatividad general. Será, por supuesto, un viaje virtual que nos conducirá a un lugar sin fronteras, el territorio de las preguntas fundamentales que se plantea la humanidad. Pero, como todo viaje, también aspira a enriquecer nuestra visión del mundo, ampliar nuestros horizontes y, al final, hacernos sentir que hemos aprendido algo. Para mí, escribir este libro ha supuesto, sin duda, todo eso.

% DOCX paragraph 45
En la ruta que propongo, haré cuanto pueda por alejarnos de las aguas traicioneras de la erudición y los tecnicismos. En cambio, nos adentraremos en mar abierto, con la intuición como referencia y la imaginación como estrella guía. Por el camino haremos algunas paradas, en parte para recuperar el aliento, dada la cantidad de conceptos nuevos que presentaré en cada capítulo, pero sobre todo para encontrar la calma necesaria para responder a preguntas sencillas, aunque no triviales, como estas:

% DOCX paragraph 47
\emph{¿Por qué una manzana cae del árbol en vez de flotar en el espacio? ¿Qué es el espaciotiempo? ¿En qué consiste su curvatura y cómo se produce? ¿Se puede curvar el tiempo? ¿Cómo funciona un agujero negro y cómo podemos «construir» uno? ¿Cómo podemos fotografiarlo si no emite luz? ¿Qué son las ondas gravitatorias y por qué son difíciles de medir?}

% DOCX paragraph 49
Como en cualquier viaje, conviene prepararse para lo que nos espera y llevar en el equipaje lo que podamos necesitar por el camino. Yo llevaré todo lo que he aprendido en treinta años dedicados al estudio de la gravedad y, en particular, de aquellos aspectos ligados de manera inseparable a la astrofísica de los agujeros negros, las estrellas de neutrones y las ondas gravitatorias. Lo aprendido durante estas décadas me ha llevado a hacer predicciones y descubrimientos; el logro más reciente llegó en abril de 2019, junto con la colaboración Event Horizon Telescope: la publicación de la primera imagen de un agujero negro supermasivo.

% DOCX paragraph 50
El equipaje de los lectores puede ser más ligero y solo debe contener dos elementos esenciales: abundante imaginación y una buena dosis de paciencia. La primera ayudará a encontrar las respuestas que «los del oficio» podemos leer fácilmente en las ecuaciones. La segunda será útil porque no todo lo que escriba resultará claro de inmediato, aunque puedo garantizar que será correcto, ni todo lo que se lea parecerá evidente, quizá ni siquiera razonable. Sin embargo, quien se arme de imaginación y paciencia puede tener la certeza de que encontrará respuesta a cada una de las preguntas anteriores y llegará a comprender el papel del espaciotiempo y la curvatura a la hora de explicar qué es realmente la gravedad, esa fuerza misteriosa que nos atrae a todos.

% DOCX paragraph 51
Comenzaremos nuestro viaje con un instinto irracional y neonatal que nos llevará hasta las orillas del más puro asombro. Allí comprenderemos qué es realmente la gravedad y cómo funcionan algunas de sus manifestaciones más extrañas, como las estrellas de neutrones y los agujeros negros.

% DOCX paragraph 52
Este será nuestro viaje, del instinto al asombro, por el revolucionario universo de Einstein.
